\scanpage{549}
\mspage{460}

$\alpha$ \deleted{\unsure{non} limite} on peut trouver un
\[
\deleted{\unsure{x}}\xi_\alpha\in E_{i_\alpha}\setminus\bigcup_{\beta<\alpha}E_{i_\beta}.
\]
L'application
\diagram{\begin{tikzcd}[ampersand replacement=\&,row sep=small,column sep=small]
\alpha\arrow[r,phantom,"\unsure{\mapstochar\relbar|\mapsto}" description]\&\deleted{\unsure{x}}\xi_\alpha\\
I_\alpha\arrow[r,phantom,"\unsure{-|\to}" description]\&E'_\infty
\end{tikzcd}}
est \deleted{\unsure{évi}}injective, donc $\operatorname{Card} E'_\infty\geq\operatorname{Card} I_{\alpha_0}$

\uline{Corollaire} 4.15.8.4. Pour \inserted{un $M$ et un $\pi_0$ donnés, pour} que le lemme-clef 4.15.8.1 soit valable, il f et s.
\noindent\makebox[0pt][r]{\smash{\rule[-10\baselineskip]{0.4pt}{10\baselineskip}}\quad}qu'il le soit \deleted{pour, \unsure{un}} dans le cas où $I=I_{\alpha_0}$, $\alpha_0$ étant le premier ordinal tel que $\operatorname{card}\alpha_0>\pi_0$.
Il revient au $=$ qu'on ne puisse trouver de système inductif $(\alpha_{ji})_{j\geq i}$ indexé par $I_{\alpha_0}$, à valeurs dans $M$, tel que pour \deleted{\unsure{j|j_0|j>}} $i<j$, on ait $\alpha_{ji}$ non inversible dans $M$.
% archive sheet number: 548
